\documentclass[12pt]{article}

\usepackage[a4paper, total={6in, 8in}]{geometry}
\usepackage{textcomp}

\begin{document}
\setlength{\parindent}{0pt}

\def \t {\textrightarrow}
\def \v {\vspace{1em}}
\def \bi {\begin{itemize}}
\def \ei {\end{itemize}}
\def \s[#1] {\section*{#1}}
\def \ss[#1] {\subsection*{#1}}
\def \sss[#1] {\subsubsection*{#1}}

\s[D'annunzio - Summa]
Ermione è muta \t non è solo lei la donna oggetto \t c'è anche Ippolita nel trionfo della morte, le vergini della roccia, la Muti \t la donna è oggettificata, visione maschilista \\
Donna viene percepito come oggetto da esplorare e abbandonare \\
Non vive per se, ma in funzione \t la donna è recepita come oggetto del desiderio (come nella pioggia del pineto, che si configura un terzo attore nell atto amoroso che è la natura \t anche lei partecipa) \\
Laddove i due si trovano, c'è uno spettatore e attore che è la natura \t la complicita amorosa dei due corpi, ma vista in presenza di uno spettatore muto \\
Le faccettature di lettura di d'annunzio sono molteplici \\
La natura è viva, ma non parla \t come la luna \t partecipa \\
D'annunzio non è solo voyerismo e scarsa considerazione della donna, ma un bisogno irrisolto che probabilmente porta a un continuo esplorare

\v

C'è la percezione di un pubblico, percepito come folla \t e la percezione della donna associata a un volto \t e la percezione di una donna particolare che chiude il cerchio: la madre \\
Il poema paradisiaco porta il ritorno alla madre, all'acqua con il liquido amniotico \\
C'è anche però l'innovazione tecnologica \t D'annunzio va in volo \t e capitalismo, sprezza la massa ma non il capitale \\
La psicologia della folla, Gustav de Bon \t la folla è irritabile, ha inclinazione a seguire un forte (che ha forte responsabilita) \t dinamica della folla è la dimanica del gruppo \\
Si vede anche nelle dinamiche giovanili \\
La folla è anche istintiva, e si adegua meccanicamente alle voci di un altro \t e l'improvvisa mutevolezza della fola = spersonalizzazione della folla \t che si vede nella dinamica dei followers \\
Omologazione della folla \t Le Bon con psicologia sociologica e sociologia psicologica

\s[Carducci]
Generazione romantica sfocia nel decadentismo \t in Pascoli si ritrovano elementi leopardiani \\
Lui è il vate della terza roma, il vate della nuova roma \t è l'ultimo dei romantici e il primo della nuova era \\
Viene visto nei suoi aspetti + esponenzialmente esposti nella politica, e visto come occasionale nelle sue scelte strategiche \t è anticlericale \\
Antipatia che trascorre tra Pascoli e Carducci

\ss[Vita]
Nasce nel 1835 in Versiglia \t in Toscana \t vive a contatto con il paesaggio Maremmano \\
Studia a Pisa, si lauera in filosofia e filologia \t ha passione sfrenata per la classicità \\
Insegnera al Ginnasio di San Miniato, ma a causa del suo essere filo-repubblicano viene allontanato \\
Vena incline al socialismo, e acceso anticlericalismo in "Inno a Satana" = ribaltamento della visione religiosa, e superamento del manzonismo \\
Pubblica "Levia gravia" (cose leggere cose pesanti, in latino \t ossimoro) \t nel titolo si vede gia amore per classicita \\
Ha passione letteraria \t che a volte sfocia nell'invettiva \\
"Giambi ed Epodi" \t ne abbiamo parlato con Dante \\
"Odi barbare" \t il contrasto bianco-nero, e anche qua il contrasto \t se una personalità del mondo classico ascoltasse le nostre liriche, le chiamerebbe barbariche \t come era barbaro chiunque non parlava latino \\
Un fatto determina la sua evoluzione di uomo e di artista \t il figlio Dante a 3 anni muore nel 1860 \t dolorosa esperienza che molto influenza la sua produzione \\
Dedica nel 71 "Pianto antico" al figlio \\
Queste liriche mettono in evidenza il suo mondo \t ma allora ha abbandonato la politica per svolta + intimistica? no \t pero ha abbandonato gli ideali + radicali e repubblicani a favore della monarchia \\
La critica parla di opportunismo \\
Il suo essere ondivago è influenzato dalla instabilità politica dell'epoca \t col fascismo ci sarà una presenza forte \\
Ondeggia politicamente, che si differenzia dal trasformismo di de pretis

\v

In tutto questo ebbe un'amante, che si trova nella lirica "Stazione di un mattino d'autunno" \t lirica che mette alla luce il progresso, e il distacco dalle persone dove un avvicinamento è possibile \t mondo digitale \\
Parla della sua donna amata mentre un treno la porta via \t il paradosso del progresso si manifesta, nel contrasto tra lontano e vicino \\
Il suo amore per Carolina Cristofori Piva \\
Nel 1890 diventa senatore del regno unito \t cosi diventa "Poeta Vate" \\
Poi riceve il premio nobel per la letteratura nel 1906, ma nel 1907 muore \\
Getta le basi per la svolta, per il verso livero, per una parabola della poesia che rompe con il romanticismo \t e che si apre a un nuovo mondo

\ss[San Martino]
Il testo tratteggia molto \t l'animo e il paesaggio \\
L'io lirico, che mostra come una pittura a tinte graduate del paesaggio autunnale (riferimento a Lavandare, e mondo autunnale di Novembre) \\
San Martino (11 Novembre) \t qui è presente la calorosita del momento attraverso il ribollire dei tini e il profumo dei vini, che creano aspetto ludico dell'atmosfera

\v

Nella prima strofa si trova la nebbia \t i colli irti creano quasi un limite \\
Piovigginando sale \t carattere fortemente iperbatico del testo \\
Maestrale è un vento freddo \t e il mare urla, personificazione \t e segno di interpunzione medio-forte (;), con urla sinestesia e onomatopea 

\v

Rottura sancita alla seconda strofa con Ma \\
Borgo \t leopardi \t le vie del borgo, ci sono le aziende vinicole \\
I tini ribollono, come i sentimenti nell'animo del poeta \\
La maturazione del vino assume la sfumatura asprigna e amarognola, che riporta al decadentismo d'annunziano \\
I tini sono come il ventre di una madre che sta generando \t e la finalita del vino è l'immagine di una coralita e leggerezza dell'animo \\
La frattura non si fa attendere dopo rallegrar: c'è un punto, segno interpunzione forte

\v

Fuoco in una casa di campagna, lo spiedo al fuoco = espressione dei piaceri della vita, di matrice quasi epicurea \t lo spiedo preannuncia ancora la coralita \\
Il cacciatore \t riporta al borgo leopardiano \t un operosita che guarda dall'uscio, che è un limite \t non vede la totalita

\v

La pensosita è tradotta nell'immagine di questi storni di uccelli \t tra le nubi rosseggianti del tramonto, come paragonato a pensieri che vagheggiano nell'aria della sera \\
Dimensione serale \t momento di introspezione, e di guardare dentro l'anima \\
Scrive la lirica a Roma \t ma la Maremma è teatro delle sue composizioni \t è il suo lougo dell'animo, ed è il mondo che lo vede protagonista di una vita diversa, che parte e ritorna sempre nello stesso punto \\
Nella lirica ci sono provocazioni coloristiche, oggettistiche e altre dimensioni del limite
\end{document}
